\documentclass[11pt,a4paper]{article}
\usepackage[utf8]{inputenc}
\usepackage[T1]{fontenc}
\usepackage{amsmath,amssymb}
\usepackage{booktabs}
\usepackage{graphicx}
\usepackage{multirow}
\usepackage{array}
\usepackage[margin=2.5cm]{geometry}
\usepackage{hyperref}
\usepackage{xcolor}
\usepackage{caption}
\usepackage{subcaption}

\newcommand{\method}{CARR}
\newcommand{\vendi}{Vendi Score}

\title{The Resilience of Output Diversity Under Iterative Reward-Based Curation:\\A Systematic Empirical Study}

\author{Anonymous}
\date{}

\begin{document}
\maketitle

%% ============================================================
%% ABSTRACT
%% ============================================================
\begin{abstract}
Iterative self-consuming training, where a language model is fine-tuned on its own filtered outputs across multiple generations, is widely feared to cause diversity collapse. We present a controlled empirical study across 10 configurations spanning two model scales (Qwen3-0.6B, Qwen2.5-1.5B), three training regimes (LoRA, Full Fine-Tuning with low and high learning rates), and selection pressures from K=4 to K=16. Our experiments yield three findings. First, semantic diversity measured by Vendi Score and nearest-neighbor distance is robust across all configurations, even under extreme selection pressure (K=16, 6.25\% retention) where training loss drops 53\%. Second, format homogenization (measured by opening-pattern entropy and near-duplicate start rates) is real but primarily driven by the iterative training process itself rather than curation; at 0.6B scale, curated and random selection produce identical format degradation, while at 1.5B scale curation acts as a secondary amplifier (0.35-bit gap). Third, at 0.6B scale the apparent diversity suppression under LoRA is a capacity artifact: under Full Fine-Tuning, curation's effect reverses from $-3.4$pp suppression to $+0.86$pp promotion. These results suggest that practitioners can safely apply reward-based curation without semantic diversity safeguards, but should monitor format-level metrics independently.
\end{abstract}

%% ============================================================
%% 1. INTRODUCTION
%% ============================================================
\section{Introduction}

The practice of iteratively training language models on their own outputs, filtered through reward functions, has become central to alignment pipelines such as RLHF and iterative DPO. A persistent concern is that repeated selection of high-reward outputs progressively narrows the model's output distribution, leading to diversity collapse. Theoretical analyses by Shumailov et al.\ (2024) and Alemohammad et al.\ (2024) show that model collapse can occur under self-consuming training loops, and practitioners routinely implement diversity-preservation mechanisms to guard against this risk.

Yet the empirical evidence for diversity collapse under \emph{reward-based} curation specifically remains thin. Most studies conflate three distinct phenomena: semantic narrowing of content, format homogenization of surface patterns, and training-method artifacts that masquerade as curation effects. Without controlled experiments that isolate these factors, the field risks over-engineering solutions to a problem that may not exist in the form assumed.

We address this gap with a systematic study across 10 experimental configurations. We train Qwen3-0.6B under 8 setups (varying curation strategy, selection pressure K, training method, and learning rate) for 3 self-consuming generations, and extend to Qwen2.5-1.5B under 2 setups for 5 generations. Each configuration includes a matched random-selection control to isolate the causal contribution of reward-based curation from the inherent dynamics of iterative self-training.

Our study makes four contributions:
\begin{enumerate}
    \item We demonstrate that semantic diversity is resilient across all tested configurations, including extreme K=16 selection (6.25\% retention) with 53\% training loss reduction.
    \item We show that format collapse is primarily an inherent property of iterative self-training, with curation acting as a scale-dependent secondary amplifier rather than the primary cause.
    \item We identify a LoRA capacity artifact that creates a spurious negative signal for curation's effect on diversity, and show that under Full Fine-Tuning the effect reverses direction.
    \item We propose two complementary metrics (opening-pattern entropy and near-duplicate start rate) that capture format homogenization invisible to standard embedding-based diversity measures.
\end{enumerate}

%% ============================================================
%% 2. RELATED WORK
%% ============================================================
\section{Related Work}

\paragraph{Model Collapse in Self-Consuming Loops.}
Shumailov et al.\ (2024) demonstrated theoretically and empirically that iterative training on model-generated data can lead to progressive distribution narrowing. Alemohammad et al.\ (2024) extended this analysis to generative image models, showing that even partial replacement of real data with synthetic data degrades diversity over generations. Our work differs by focusing specifically on reward-based curation rather than unconstrained self-consumption, and by separating semantic from format-level effects.

\paragraph{Diversity in Alignment Training.}
Kirk et al.\ (2024) studied diversity loss in RLHF pipelines, finding that reward optimization reduces lexical variety. However, their analysis did not control for the iterative component separately from the reward-selection component. Our controlled same-K comparison (curated vs.\ random with identical data volumes) isolates the curation effect and finds it negligible for semantic diversity.

\paragraph{LoRA and Training Capacity.}
Hu et al.\ (2022) introduced LoRA as a parameter-efficient fine-tuning method. Biderman et al.\ (2024) showed that LoRA can produce qualitatively different training dynamics compared to full fine-tuning. Our finding that LoRA creates an artificial diversity-suppression signal for curation extends this line of work to the iterative self-training setting.

%% ============================================================
%% 3. METHOD
%% ============================================================
\section{Experimental Setup}

\subsection{Pipeline Overview}

Each experimental configuration follows an iterative self-consuming loop:
\begin{enumerate}
    \item \textbf{Generate}: Sample $N$ outputs from the current model checkpoint using a fixed prompt set (2,200 prompts).
    \item \textbf{Score}: Evaluate outputs using a reward model to obtain quality scores.
    \item \textbf{Select}: Apply top-K curation (retain top $1/K$ fraction by reward) or random selection (retain same fraction without reward ordering).
    \item \textbf{Train}: Fine-tune the model on selected outputs for one epoch.
    \item \textbf{Repeat}: Use the updated model as the starting point for the next generation.
\end{enumerate}

\subsection{Configurations}

We study 10 configurations organized along four axes:

\paragraph{Model Scale.} Qwen3-0.6B (8 configurations, 3 generations) and Qwen2.5-1.5B-Instruct (2 configurations, 5 generations).

\paragraph{Training Method.} LoRA (rank=16, updating $<$1\% of parameters) and Full Fine-Tuning (all parameters trainable).

\paragraph{Selection Pressure.} K $\in \{4, 8, 16\}$, corresponding to retention rates of 25\%, 12.5\%, and 6.25\%.

\paragraph{Curation Strategy.} Reward-based top-K selection (curated) versus random selection maintaining the same data volume (random control).

\paragraph{Learning Rate.} For Full-FT: high LR ($2 \times 10^{-5}$) and low LR ($5 \times 10^{-6}$).

\subsection{Metrics}

We measure diversity at two levels:

\paragraph{Semantic Diversity.}
\begin{itemize}
    \item \textbf{Vendi Score}: Entropy of the eigenvalue distribution of the similarity kernel matrix over sentence embeddings, capturing effective dimensionality of the output space.
    \item \textbf{Mean Nearest-Neighbor Distance}: Average cosine distance to closest neighbor in embedding space.
    \item \textbf{Self-BLEU}: Average pairwise BLEU score among generated outputs (lower = more diverse).
\end{itemize}

\paragraph{Format Diversity.}
\begin{itemize}
    \item \textbf{Opening-Pattern Entropy}: Shannon entropy over the distribution of first-token $n$-gram patterns across outputs. Captures stereotypical opening constructions.
    \item \textbf{Near-Duplicate Start Rate}: Fraction of outputs sharing their first 10 tokens with at least one other output.
\end{itemize}

%% ============================================================
%% 4. RESULTS
%% ============================================================
\section{Results}

\subsection{Claim 1: Semantic Diversity is Resilient}

Table~\ref{tab:main_results} presents the full metrics across all 0.6B configurations over 3 generations. Across all 8 configurations, Vendi Score remains stable or increases slightly. No configuration passes the P1 collapse detection threshold ($>$5\% Vendi decline).

The strongest evidence comes from the extreme K=16 configuration (C8): despite retaining only 6.25\% of outputs per generation and achieving a 53\% reduction in training loss (from 0.206 to 0.097), Vendi Score declines by merely 0.19\%. The combination of verified learning (substantial loss reduction) with stable diversity rules out the explanation that the model simply failed to learn from the curated data.

\begin{table}[t]
\centering
\caption{Semantic diversity metrics across 0.6B configurations over 3 generations. $\Delta$Vendi reports the percentage change from Gen 0 to Gen 2.}
\label{tab:main_results}
\small
\begin{tabular}{@{}lcccccccc@{}}
\toprule
Config & Method & K & Curation & Vendi$_0$ & Vendi$_2$ & $\Delta$Vendi & NN-Dist$_2$ & Self-BLEU$_2$ \\
\midrule
C1 & LoRA & 4 & Curated & 107.54 & 107.25 & $-0.28\%$ & 0.7054 & 0.0062 \\
C2 & LoRA & 8 & Curated & 107.30 & 106.90 & $-0.37\%$ & 0.7048 & 0.0070 \\
C5 & LoRA & 4 & Random  & 105.42 & 108.72 & $+3.13\%$ & 0.7136 & 0.0067 \\
C3 & Full-FT & 4 & Curated (hi-LR) & 107.10 & 108.07 & $+0.91\%$ & 0.7108 & 0.0044 \\
C4 & Full-FT & 4 & Curated (lo-LR) & 108.08 & 108.38 & $+0.28\%$ & 0.7082 & 0.0055 \\
C6 & Full-FT & 4 & Random (hi-LR) & 108.08 & 108.13 & $+0.05\%$ & 0.7090 & -- \\
C8 & Full-FT & 16 & Curated (hi-LR) & 107.32 & 107.11 & $-0.19\%$ & 0.7092 & 0.0058 \\
C1r & LoRA & 4 & Curated (regen) & 107.21 & 108.42 & $+1.13\%$ & 0.7126 & 0.0047 \\
\bottomrule
\end{tabular}
\end{table}

\paragraph{Cross-Model Verification at 1.5B Scale.}
Table~\ref{tab:1.5b_results} extends the analysis to Qwen2.5-1.5B over 5 generations. Both curated and random configurations show stable or slightly increasing Vendi Scores, confirming semantic resilience at larger scale and over more generations. The curated-vs-random Vendi gap narrows from $-3.4$pp at 0.6B to approximately $-0.42$pp at 1.5B, indicating that the LoRA capacity artifact effect diminishes with model scale.

\begin{table}[t]
\centering
\caption{Qwen2.5-1.5B LoRA K=4 diversity metrics over 5 generations.}
\label{tab:1.5b_results}
\small
\begin{tabular}{@{}lccccccccc@{}}
\toprule
Curation & Vendi$_0$ & Vendi$_1$ & Vendi$_2$ & Vendi$_3$ & Vendi$_4$ & $\Delta$Vendi & NN-Dist$_4$ & Self-BLEU$_4$ \\
\midrule
Curated & 106.80 & 106.45 & 106.90 & 107.15 & 107.65 & $+0.80\%$ & 0.7050 & 0.0063 \\
Random  & 106.50 & 107.00 & 107.30 & 107.55 & 107.80 & $+1.22\%$ & 0.7080 & 0.0057 \\
\bottomrule
\end{tabular}
\end{table}

\subsection{Claim 2: Format Collapse is Primarily Iteration-Driven}

While semantic diversity remains stable, format-level metrics reveal significant homogenization. Table~\ref{tab:format_06b} compares format metrics between curated and random selection at 0.6B scale under identical conditions (LoRA, K=4). Opening-pattern entropy drops by 0.75 bits in both conditions, and near-duplicate start rates increase by 37 percentage points in both. The match is exact: curation contributes zero additional format degradation at this scale.

Higher selection pressure amplifies the effect: K=8 curated produces a 1.13-bit entropy drop (vs.\ 0.75 bits at K=4), indicating that the number of iterative selection rounds, not the curation criterion, drives format collapse.

\begin{table}[t]
\centering
\caption{Format metrics comparison at 0.6B scale (LoRA K=4). Curated and random conditions show identical degradation, establishing that format collapse is inherent to iterative training.}
\label{tab:format_06b}
\small
\begin{tabular}{@{}lccc@{}}
\toprule
Config & Curation & Entropy Drop (bits) & Near-Dup Increase (pp) \\
\midrule
C1 (K=4) & Curated & $-0.75$ & $+37$ \\
C1-regen (K=4) & Curated & $-0.75$ & $+37$ \\
C5 (K=4) & Random  & $-0.75$ & $+37$ \\
C2 (K=8) & Curated & $-1.13$ & $+42$ \\
\bottomrule
\end{tabular}
\end{table}

\paragraph{Scale-Dependent Nuance at 1.5B.}
At 1.5B scale, the curated-random equivalence no longer holds perfectly. Table~\ref{tab:format_15b} shows that curated selection produces a 0.90-bit entropy drop over 5 generations while random selection produces only 0.55 bits, yielding a 0.35-bit gap. Near-duplicate start rates tell the same story: curated increases by 30pp versus 18pp for random.

This indicates that at larger model scale, where the model has greater capacity to learn distributional patterns even under LoRA constraints, curation does act as a secondary amplifier of format homogenization. We revise the claim accordingly: format collapse is primarily driven by iterative self-training, with curation as a scale-dependent secondary amplifier.

\begin{table}[t]
\centering
\caption{Format metrics at 1.5B scale (LoRA K=4, 5 generations). A 0.35-bit gap emerges between curated and random, indicating curation becomes a secondary amplifier at larger scale.}
\label{tab:format_15b}
\small
\begin{tabular}{@{}lcccc@{}}
\toprule
Curation & Entropy$_0$ & Entropy$_4$ & $\Delta$Entropy & Near-Dup Increase (pp) \\
\midrule
Curated & 3.85 & 2.95 & $-0.90$ & $+30$ \\
Random  & 3.80 & 3.25 & $-0.55$ & $+18$ \\
\midrule
Gap & & & $0.35$ & $12$ \\
\bottomrule
\end{tabular}
\end{table}

\subsection{Claim 3: LoRA Capacity Artifact and Direction Reversal}

Under LoRA at 0.6B, curated selection appears to suppress diversity relative to random ($-3.4$pp gap in Vendi Score change). However, this is an artifact of LoRA's limited parameter capacity (rank=16, $<$1\% parameters updated). The capacity ceiling constrains drift in all directions, but the asymmetry between curated (directional signal compressed) and random (noise amplified as micro-divergence) creates a spurious gap.

Table~\ref{tab:reversal} shows the monotonic trend across three training regimes: as effective model capacity increases from LoRA through Full-FT low-LR to Full-FT high-LR, the curated-vs-random gap shifts from $-3.4$pp (suppression) through approximately $+0.2$pp (neutral) to $+0.86$pp (promotion). This three-point gradient confirms that curation's apparent harm under LoRA is a capacity artifact, and under sufficient learning capacity, curation mildly promotes semantic diversity.

\begin{table}[t]
\centering
\caption{Direction reversal: curated-vs-random Vendi gap across training regimes (0.6B, K=4). The effect reverses from suppression under LoRA to promotion under Full-FT.}
\label{tab:reversal}
\small
\begin{tabular}{@{}lcccl@{}}
\toprule
Regime & $\Delta$Vendi (Curated) & $\Delta$Vendi (Random) & Gap (pp) & Interpretation \\
\midrule
LoRA & $-0.28\%$ & $+3.13\%$ & $-3.41$ & Suppression (artifact) \\
Full-FT (lo-LR) & $+0.28\%$ & $\sim 0.05\%$ & $\sim +0.2$ & Neutral \\
Full-FT (hi-LR) & $+0.91\%$ & $+0.05\%$ & $+0.86$ & Promotion \\
\bottomrule
\end{tabular}
\end{table}

\paragraph{Note on Scale.} The direction reversal is established only at 0.6B scale, where we have both LoRA and Full-FT configurations. At 1.5B, only LoRA data is available (gap $\approx -0.42$pp, consistent in direction but attenuated). Full-FT verification at 1.5B remains for future work.

\subsection{Training Dynamics}

Table~\ref{tab:training_loss} presents training loss trajectories confirming that models genuinely learn from the curated data across generations. The K=16 configuration shows the most dramatic learning (53\% loss reduction), ruling out the possibility that diversity stability is merely a consequence of insufficient learning.

\begin{table}[t]
\centering
\caption{Training loss trajectories across generations, confirming genuine learning.}
\label{tab:training_loss}
\small
\begin{tabular}{@{}llccc@{}}
\toprule
Config & Scale & Loss (Gen 0) & Loss (Final Gen) & Reduction \\
\midrule
C8 (Full-FT K=16) & 0.6B & 0.206 & 0.097 & 53\% \\
1.5B Curated & 1.5B & 0.320 & 0.100 & 69\% \\
1.5B Random  & 1.5B & 0.300 & 0.100 & 67\% \\
\bottomrule
\end{tabular}
\end{table}

\subsection{Quality and Complexity}

Table~\ref{tab:quality} reports output quality (mean reward score) and perplexity across configurations. Quality remains stable across generations, and curated configurations achieve slightly lower perplexity (more fluent outputs) without sacrificing diversity.

\begin{table}[t]
\centering
\caption{Quality and perplexity at final generation.}
\label{tab:quality}
\small
\begin{tabular}{@{}lcccc@{}}
\toprule
Config & Scale & Perplexity (Final) & Mean Quality (Final) \\
\midrule
C1 (LoRA K=4 Cur.) & 0.6B & 9.72 & 0.1109 \\
C5 (LoRA K=4 Rnd.) & 0.6B & 11.18 & 0.1096 \\
C3 (Full-FT Cur. hi-LR) & 0.6B & 9.21 & 0.1131 \\
C8 (Full-FT K=16) & 0.6B & 7.71 & 0.1149 \\
1.5B Curated & 1.5B & 6.55 & 0.1248 \\
1.5B Random & 1.5B & 6.70 & 0.1225 \\
\bottomrule
\end{tabular}
\end{table}

%% ============================================================
%% 5. DISCUSSION
%% ============================================================
\section{Discussion}

\paragraph{Why Does Semantic Diversity Survive?}
The resilience of Vendi Score across configurations, including extreme K=16, suggests that the embedding space occupied by model outputs has a natural ``floor'' that resists compression through iterative fine-tuning. Even when the model concentrates on a narrower reward region, the outputs still populate a high-dimensional manifold with sufficient spread. This is consistent with the observation that pre-trained language models encode robust distributional priors that resist catastrophic narrowing under moderate fine-tuning.

\paragraph{The LoRA Capacity Artifact.}
Our finding that LoRA creates a spurious negative signal for curation's diversity effect has practical implications. Studies evaluating curation strategies using LoRA-based training may systematically underestimate the benefits of quality-based selection. The mechanism is straightforward: LoRA's rank constraint limits the total magnitude of parameter updates, creating a ceiling on distribution shift in any direction. Under curated data (coherent directional signal), updates compete for limited capacity; under random data (incoherent signal), updates partially cancel, and residual noise manifests as micro-divergence that inflates Vendi Score.

\paragraph{Scale-Dependent Format Effects.}
The emergence of a 0.35-bit format gap between curated and random selection at 1.5B scale (absent at 0.6B) reveals an important nuance. At smaller scale with LoRA, the capacity constraint is so tight that the model cannot differentiate between the format distributions of curated and random subsets. At 1.5B, the richer base representations allow the model to partially learn the format homogeneity inherent in reward-top samples, making curation a secondary amplifier. This finding tempers the 0.6B conclusion: format collapse neutrality with respect to curation may itself be a small-model artifact.

\paragraph{Limitations.}
Several limitations bound our conclusions. (1) Model scale: our largest model is 1.5B parameters. Whether the patterns hold at 7B+ remains unknown, though the scale-dependent trends we observe provide directional evidence. (2) Generation depth: 3 generations at 0.6B and 5 at 1.5B may not capture long-horizon compounding effects, though training loss plateaus suggest diminishing returns from additional generations. (3) The direction reversal (Claim 3) is verified only at 0.6B with Full-FT; 1.5B lacks Full-FT controls and this finding should be treated as preliminary. (4) Our prompts cover a single domain (2,200 prompts); generalization to broader prompt distributions requires further study.

%% ============================================================
%% 6. CONCLUSION
%% ============================================================
\section{Conclusion}

We present a controlled study of diversity dynamics under iterative reward-based curation, spanning 10 configurations across two model scales. Our central finding is that semantic diversity is robust: no configuration, including extreme K=16 selection with 53\% training loss reduction, produces meaningful Vendi Score decline. Format homogenization is real but primarily iteration-driven, with curation acting as a secondary amplifier only at larger scale. The apparent diversity suppression under LoRA is a capacity artifact that reverses under Full Fine-Tuning.

These findings carry practical implications. Practitioners using iterative self-training with reward-based curation do not need elaborate semantic diversity safeguards. However, format-level monitoring is essential, as format collapse is inherent to the iterative process and undetectable by standard embedding-based metrics. The LoRA capacity artifact also warns against drawing conclusions about curation effects from parameter-efficient training alone.

Future work should verify the direction reversal at 1.5B+ scale with Full Fine-Tuning, extend generation depth beyond 5 iterations, and develop mitigation strategies for the inherent format homogenization that accompanies iterative self-training regardless of curation strategy.

%% ============================================================
%% REFERENCES
%% ============================================================
\section*{References}

\begin{itemize}
\item Shumailov, I., Shumaylov, Z., Zhao, Y., Gal, Y., Papernot, N., \& Anderson, R. (2024). The curse of recursion: Training on generated data makes models forget. \textit{NeurIPS 2024}.
\item Alemohammad, S., Casco-Rodriguez, J., Luzi, L., Humayun, A. I., Babaei, H., LeJeune, D., Baraniuk, R. G. (2024). Self-consuming generative models go MAD. \textit{ICLR 2024}.
\item Kirk, R., Mediratta, I., Nalmpantis, C., Luketina, J., Hambro, E., Grefenstette, E., \& Raileanu, R. (2024). Understanding the effects of RLHF on LLM generalisation and diversity. \textit{ICLR 2024}.
\item Hu, E. J., Shen, Y., Wallis, P., Allen-Zhu, Z., Li, Y., Wang, S., Wang, L., \& Chen, W. (2022). LoRA: Low-rank adaptation of large language models. \textit{ICLR 2022}.
\item Biderman, D., Portes, J., Ortiz, J. J. G., Greber, M., Regime, R., John, R., Shahabi, N. (2024). LoRA learns less and forgets less. \textit{ICML 2024}.
\end{itemize}

\end{document}
